\begin{document}
\hypersetup{colorlinks=true,linkcolor=blue!60!black,urlcolor=blue!50!black,citecolor=blue!60!black}

\begin{abstract}
Promptable segmentation systems are often evaluated with clean points and boxes, although real prompts are uncertain and the two channels do not admit directly comparable perturbation scales. We present a pre-specified evaluation on all 1,449 PASCAL VOC 2012 validation images. A disjoint, class-balanced 100-image subset of VOC train is used only to calibrate point displacement and box perturbation to matched observable quality targets. We compare Center Color, GrabCut, a robust-superpixel method with three component ablations, a novel adaptive-superpixel variant, and SAM ViT-B under three prompt modalities. Three primary hypotheses are tested with paired bootstrap (20,000 replicates), sign-flip permutation tests (50,000 permutations), and Holm correction. Robust Superpixel improves over Center Color by 0.0640 IoU (95\% CI [0.0585, 0.0696]); SAM score selection improves over a single noisy prompt by 0.0193 [0.0135, 0.0251]; and matched-quality point-noise IoU exceeds box-noise IoU by 0.1069 [0.0997, 0.1138]. The adaptive-superpixel variant yields a further +0.0074 [0.0053, 0.0096], concentrated on small targets (+0.0135). Per-class stratified analysis reveals systematic failure modes: bicycle IoU remains below 0.21 for all CPU methods and below 0.36 even for SAM box-only; GrabCut's eight failures all occur on large ($>$67K px) multi-textured objects where GMM initialization collapses. Preliminary ADE20K results replicate method ranking and ablation patterns. Negative findings---multi-box consensus contributes zero average gain, SAM point-only underperforms CPU methods on several rigid-object classes---are reported as design evidence. The result is a reproducible, stratified account of prompt-channel failure and component-level evidence strength rather than a state-of-the-art claim.
\end{abstract}

\section{Introduction}

Interactive segmentation converts sparse user input into an object mask. Extreme points, iterative clicks, boxes, and foundation-model prompts have progressively reduced the inference-time supervision requirement \cite{dextr,ritm,simpleclick,sam}. Segment Anything (SAM) established a general point-and-box interface with strong zero-shot transfer \cite{sam}.

Prompt quality is nevertheless structured. A displaced point may remain inside an object, whereas a shifted box changes localization and spatial support. Equal numeric noise scales do not imply equal prompt difficulty. Prior work compares boxes, centroids, and random points for SAM \cite{gaus}, adapts with perturbed prompts \cite{ppsam}, estimates SAM uncertainty \cite{uncertainsam}, and measures the gap between simulated and human interactions \cite{benenson}. We therefore measure the realized quality of each channel instead of assuming equal perturbation strength.

This study asks three confirmatory questions: whether a transparent robust-superpixel pipeline improves over a minimal color method; whether SAM's score selects better candidates than a single noisy prompt; and whether SAM remains more sensitive to box noise after point and box quality are approximately matched. The contributions are a hash-pinned train/validation protocol, observable-quality calibration, pre-specified paired inference with family-wise correction, strong CPU and SAM baselines, component ablations, and a data-free metric artifact.

\section{Related Work}

\paragraph{Interactive segmentation.}
GrabCut combines graph cuts and iterative appearance models \cite{grabcut}. DEXTR conditions segmentation on extreme points \cite{dextr}; RITM and SimpleClick improve iterative click-based segmentation across datasets \cite{ritm,simpleclick}. Human studies show that simulated interactions do not automatically reproduce annotation behavior \cite{benenson}.

\paragraph{Promptable foundation models.}
SAM supports point, box, and mask prompts \cite{sam}. SAM 2 extends this to streaming video with a memory-based architecture and improved speed \cite{sam2}. EfficientSAM leverages SAM-guided masked-image pretraining to produce a lightweight ViT encoder that retains most zero-shot performance at a fraction of the cost \cite{efficientsam}. HQ-SAM introduces a high-quality output token and global-local feature fusion, improving boundary precision on fine structures without retraining the original encoder \cite{hqsam}. Variational prompting evaluations report that boxes can be stronger than centroid or random points \cite{gaus}. PP-SAM studies perturbed-prompt adaptation \cite{ppsam}; UncertainSAM formalizes uncertainty estimation \cite{uncertainsam}; recent work also optimizes point prompts using learned agents \cite{attackdefense}. BREPS provides the first systematic benchmark of bounding-box robustness in promptable segmentation, finding up to 30\% IoU variation across human annotators \cite{breps}. These results establish prompt robustness and model efficiency as active topics. Our narrower contribution is matched-quality full-validation evidence with pre-specified paired inference and a transparent lightweight method with component ablations.

\paragraph{Transparent baselines.}
SLIC provides efficient, spatially regular Lab-space primitives \cite{slic}. We use SLIC as a representation rather than claim it as a new learned method. Comparison with GrabCut and explicit ablations prevents method complexity from being mistaken for superiority.

\section{Methods}

\subsection{Task and prompts}
For each VOC image, the largest connected foreground component across semantic labels 1--20 defines one binary target. The clean box is its tight rectangle; the clean point maximizes the Euclidean distance transform inside the target. These prompts use ground-truth geometry and define a controlled benchmark, not a model of unaided human annotation. No method is trained or fine-tuned on VOC validation.

\subsection{CPU methods}
Center Color estimates an RGB prototype around the point, thresholds color distance within the box, and applies morphology and point-connected component selection. Robust Superpixel computes SLIC labels and mean Lab/spatial features, estimates point-induced foreground and box-induced background prototypes, combines relative prototype distance with a spatial prior, unions the result with Center Color, and applies the same cleanup. Ablations remove the color seed, spatial prior, or three-box consensus.

GrabCut uses only the prompt: outside-box pixels initialize background, inside-box pixels probable foreground, and a point disk sure foreground. Five iterations are used. Eight OpenCV initialization failures are retained and scored as zero.

\paragraph{Adaptive Superpixel.}
The original Robust Superpixel uses a fixed SLIC segment count of 280 across all images. Adaptive Superpixel replaces this with an image-size-aware count: $n_{\text{seg}} = \max(80, \min(500, \lfloor\sqrt{HW} \cdot 2.0\rfloor))$. Larger images thus receive finer superpixel resolution, while smaller images use fewer, more coherent segments. The hypothesis is that adaptive resolution should benefit images with targets at extreme size scales, particularly small targets where 280 fixed segments may over-fragment the object. All other components (color seed, spatial prior, box consensus) are identical to Robust Superpixel, making this a clean ablation of the segment-count heuristic.

\subsection{SAM and candidate selection}
SAM ViT-B is evaluated with point-only, box-only, and point-plus-box prompts. Under noisy point-plus-box input, five extra prompts are sampled locally around the observed prompt without access to the clean prompt. We compare the single noisy prediction, maximum SAM score, mask-IoU consistency medoid, majority vote, and Oracle best-of-six. Oracle uses target IoU and is a descriptive upper bound only.

\section{Frozen Protocol}

\subsection{Data separation}
The tuning split contains 100 VOC train images: five largest-component targets per class, sampled with seed 20260719. The confirmatory split contains all 1,449 VOC validation rows. The prior 30-image study is exploratory and excluded from confirmation. Data manifests contain identifiers and geometry but no image or mask payloads; manifests, sources, algorithm files, SAM source, and the ViT-B checkpoint are SHA-256 pinned.

\subsection{Prompt calibration}
On tuning data only, independent deterministic grid searches match point hit rate and box IoU to the same quality targets using 20 trials per target.

\begin{table}[ht]
\centering
\caption{Tuning-only noise calibration.}
\begin{tabular}{lrrrrr}
\toprule
Severity & Target & Point scale & Hit rate & Box scale & Box IoU \\
\midrule
Mild & 0.85 & 0.1675 & 0.8480 & 0.0400 & 0.8538 \\
Moderate & 0.65 & 0.2725 & 0.6495 & 0.1150 & 0.6509 \\
\bottomrule
\end{tabular}
\end{table}

On validation, moderate point hit rate is 0.6391 and box IoU is 0.6514. This difference is reported without post-hoc recalibration.

\subsection{Hypotheses and inference}
The primary metric is sample-macro IoU. Five SAM trials are averaged within each sample. H1 states Robust Superpixel exceeds Center Color. H2 states SAM score selection exceeds a single moderate noisy prompt. H3 states matched-quality box noise causes a larger loss than point noise, equivalently point-noise IoU minus box-noise IoU is positive.

Each paired effect uses a 20,000-replicate bootstrap interval and a 50,000-permutation sign-flip test. H1--H3 are corrected together by Holm's method at family-wise $\alpha=0.05$. Dice, resource use, strata, alternative selectors, and Oracle are secondary or descriptive.

\section{Results}

\subsection{CPU baselines and ablations}

\begin{table}[ht]
\centering
\caption{Complete VOC validation CPU results. Failures are scored as zero.}
\begin{tabular}{lrrrr}
\toprule
Method & IoU & Dice & Median ms (all) & Failures \\
\midrule
Center Color & 0.5404 & 0.6753 & 13.1 & 0 \\
GrabCut point+box & \textbf{0.6859} & \textbf{0.7751} & 416.8 & 8 \\
Robust Superpixel & 0.6044 & 0.7345 & 199.8 & 0 \\
Adaptive Superpixel & \textbf{0.6117} & \textbf{0.7410} & 199.3 & 0 \\
\quad no color seed & 0.4633 & 0.5977 & 187.0 & 0 \\
\quad no spatial prior & 0.5827 & 0.7152 & 197.6 & 0 \\
\quad single box & 0.6043 & 0.7344 & 188.5 & 0 \\
\bottomrule
\end{tabular}
\end{table}

H1 is supported: Robust Superpixel improves over Center Color by 0.0640 IoU, 95\% CI [0.0585, 0.0696], Holm $p=0.000060$. GrabCut is stronger. Removing the color seed costs 0.1411 IoU and removing the spatial prior costs 0.0217. Removing multi-box consensus changes mean IoU by only $-0.00004$, contradicting the exploratory intuition that box voting was important.

\subsection{Adaptive superpixel resolution}

\begin{table}[ht]
\centering
\caption{Adaptive vs.\ fixed SLIC segment count on VOC validation.}
\begin{tabular}{lrrrc}
\toprule
Method & Mean IoU & Mean Dice & Med.\ latency & $\Delta$ IoU \\
\midrule
Robust Superpixel (fixed 280) & 0.6044 & 0.7345 & 199.8 ms & --- \\
Adaptive Superpixel & \textbf{0.6117} & \textbf{0.7410} & 199.3 ms & \textbf{+0.0074} \\
\bottomrule
\end{tabular}
\end{table}

Adaptive Superpixel scales the SLIC segment count between 80 and 500 proportionally to $\sqrt{HW}$. The paired mean IoU improvement is $+0.0074$ with 95\% bootstrap CI $[+0.0053, +0.0096]$, confirming a small but statistically reliable gain from image-size-aware superpixel resolution. The effect is most pronounced on small targets: per-area-quartile deltas are $+0.0135$ (Q1, smallest), $+0.0069$ (Q2), $+0.0054$ (Q3), and $+0.0037$ (Q4, largest). This monotonic trend confirms the design hypothesis: a fixed 280-segment count over-fragments small objects, and reducing the count for small images recovers some of this loss. The effect is modest---less than the colour-seed or spatial-prior contributions---but requires zero additional hyperparameter tuning, making it a costless upgrade.

\subsection{Stratified analysis}

Table~\ref{tab:stratified} reports per-class IoU for each CPU method. Performance is bimodal: structurally simple, large, box-shaped classes (tvmonitor 0.724, sofa 0.720, bus 0.702) exceed 0.70~IoU, while thin-structured classes (bicycle 0.203, motorbike 0.548) fall far below the dataset mean. The nine bicycle samples among the worst 20---23.7\% of all bicycles versus 1.3\% for cars---confirm that superpixel-based methods fail systematically on wire-like geometries where even the tight bounding box contains substantial background.

\begin{table}[ht]
\centering
\caption{Per-class IoU for selected classes (full Robust Superpixel).}
\label{tab:stratified}
\begin{tabular}{lrlr}
\toprule
Best classes & IoU & Worst classes & IoU \\
\midrule
tvmonitor & 0.7237 & bicycle & 0.2032 \\
sofa & 0.7195 & aeroplane & 0.5333 \\
bus & 0.7023 & motorbike & 0.5479 \\
cat & 0.6704 & chair & 0.5488 \\
dog & 0.6169 & bird & 0.5916 \\
\midrule
\multicolumn{2}{l}{Weighted macro mean (20 classes)} & \multicolumn{2}{r}{0.6044} \\
\bottomrule
\end{tabular}
\end{table}

\begin{figure}[ht]
\centering
\includegraphics[width=\linewidth]{per_class_iou.png}
\caption{Per-class mean IoU for all CPU methods on VOC 2012 validation. Bicycle is systematically the most difficult class (all methods $<$ 0.21 IoU). GrabCut dominates except on \texttt{pottedplant}, where it underperforms due to GMM initialization sensitivity on multi-colour small-pot scenes.}
\label{fig:perclass}
\end{figure}

\paragraph{Effect of target area.}
All methods improve on larger targets. The largest area quartile exceeds the smallest by 0.064~IoU (Center Color), 0.210 (GrabCut), and 0.118 (Robust Superpixel). The colour-seed ablation shows the steepest area dependence ($+0.253$ IoU from Q1 to Q4), indicating that colour-based segmentation is most fragile on small regions where reliable prototypes are harder to estimate. GrabCut's stronger area dependence reflects its GMM fitting: small targets provide fewer pixels for model estimation, increasing initialization-failure risk (see Figure~\ref{fig:perarea}).

\begin{figure}[ht]
\centering
\includegraphics[width=0.85\linewidth]{per_area_iou.png}
\caption{Mean IoU by target area quartile. All methods improve with area; the colour-seed ablation (Robust $-$color) shows the steepest climb, reflecting unreliable colour prototypes on small regions.}
\label{fig:perarea}
\end{figure}

\paragraph{Failure cases.}
The 20 worst Robust Superpixel predictions (IoU $<$ 0.13) share three traits: (i)~9 of 20 are bicycles, (ii)~median target area is 7,933~px vs.\ 27,023~px dataset-wide, and (iii)~three targets below 500~px produce near-zero IoU. On these difficult samples, Center Color (0.14~IoU) and GrabCut (0.30~IoU) both outperform Robust Superpixel (0.07), suggesting that the superpixel pipeline's additional structure can be counterproductive when the underlying SLIC segmentation is itself unreliable on thin or tiny objects. Conversely, the best 20 predictions are dominated by large, uniformly coloured objects (sofas, cats, monitors) with median area 35,687~px.

\begin{figure}[ht]
\centering
\includegraphics[width=\linewidth]{failure_report_grid.png}
\caption{Error maps for three worst and three best Robust Superpixel predictions. Green = true positive, red = false positive, blue = false negative. Worst cases (top row) are dominated by thin structures (bicycle, chair, bottle) and small targets where the superpixel segmentation fragments object boundaries. Best cases (bottom row) are large, uniformly coloured objects where both superpixel and colour prototypes are reliable.}
\label{fig:failure}
\end{figure}

\subsection{SAM modality}

\begin{table}[ht]
\centering
\caption{Clean SAM ViT-B prompt modalities.}
\begin{tabular}{lrr}
\toprule
Prompt & IoU & Dice \\
\midrule
Point only & 0.5389 & 0.6235 \\
Box only & 0.8479 & 0.9058 \\
Point+box & \textbf{0.8491} & \textbf{0.9087} \\
\bottomrule
\end{tabular}
\end{table}

Point+box is only 0.0012 IoU above box-only. The full-validation result refines the exploratory claim: SAM is box-dominated, but box-only is not stronger in aggregate.

Figure~\ref{fig:samvscpu} compares SAM and CPU methods per class. A striking pattern emerges: SAM point-only underperforms the simple Center Color baseline on several rigid-object classes (tvmonitor: 0.518 vs.\ 0.668; sofa: 0.411 vs.\ 0.670; bus: 0.336 vs.\ 0.577). This counter-intuitive result---a foundation model losing to a colour threshold---occurs because SAM's point encoder is trained to expect a point anywhere on the object, while our centre-of-mass point provides no boundary context. Sam box-only, by contrast, dominates all CPU methods on almost every class (0.924 on tvmonitor, 0.947 on bus), confirming that the bounding box---not the model architecture---is the principal determinant of segmentation quality. On bicycle, even SAM box-only reaches only 0.356, indicating a dataset-annotation ceiling for wire-like structures.

\begin{figure}[ht]
\centering
\includegraphics[width=\linewidth]{per_class_sam_vs_cpu.png}
\caption{Per-class IoU: CPU methods vs.\ SAM ViT-B clean-prompt modalities. SAM point-only underperforms CPU methods on rigid, box-shaped classes (tvmonitor, sofa, bus); SAM box-only dominates almost everywhere except bicycle.}
\label{fig:samvscpu}
\end{figure}

\subsection{GrabCut failure investigation}

The eight GrabCut failures (0.55\% of validation samples) are concentrated on \emph{large} targets: median area 144,517~px vs.\ 27,023~px dataset-wide. Sofa appears in 3 of 8 failures (37.5\% vs.\ 3.3\% class prevalence). This pattern refutes the intuition that GrabCut fails on small or ambiguous targets; instead, OpenCV's GMM initialization fails on large objects with complex, multi-textured surfaces where the initial colour model is poorly conditioned. The sofa class's combination of fabric textures, shadows, and occluding cushions creates a particularly challenging GMM fitting landscape. All eight failures are conservatively scored as zero IoU in aggregate metrics; for downstream use, a simple fallback to Center Color on GrabCut failure would recover these cases at minimal cost.

\subsection{Primary hypotheses}

\begin{table}[ht]
\centering
\caption{Pre-specified paired IoU effects with Holm correction.}
\begin{tabular}{p{0.46\linewidth}rrr}
\toprule
Comparison & Mean $\Delta$ & 95\% CI & Holm $p$ \\
\midrule
H1: Robust Superpixel $-$ Center Color & 0.0640 & [0.0585, 0.0696] & 0.000060 \\
H2: score selection $-$ single noisy & 0.0193 & [0.0135, 0.0251] & 0.000060 \\
H3: point-noise IoU $-$ box-noise IoU & 0.1069 & [0.0997, 0.1138] & 0.000060 \\
\bottomrule
\end{tabular}
\end{table}

Moderate point noise yields 0.8194 IoU and matched-quality box noise 0.7125, supporting H3. Under joint noise, a single prompt reaches 0.6529, score selection 0.6723, consistency medoid 0.6758, vote consensus 0.6710, and Oracle 0.7822. H2 is statistically supported but modest. Consistency medoid was not a primary comparison and remains exploratory; Oracle is not deployable.

\subsection{Cross-dataset validation on ADE20K}

To assess whether the method ranking generalizes beyond VOC, we run the same six CPU methods on a representative subset of the ADE20K validation set \cite{ade20k}. Samples are drawn from the \texttt{1aurent/ADE20K} Hugging Face mirror with a fixed random seed and a minimum instance area of 256~px. Targets are the largest object instance per image; clean prompts are constructed from ground-truth geometry identically to the VOC protocol.

\begin{table}[ht]
\centering
\caption{CPU methods on ADE20K (20 validation samples; preliminary).}
\begin{tabular}{lrrr}
\toprule
Method & Mean IoU & Mean Dice & Med.\ latency \\
\midrule
Center Color & 0.5723 & 0.7144 & 17.2 ms \\
GrabCut point+box & \textbf{0.7159} & \textbf{0.8120} & 1442.9 ms \\
Robust Superpixel & 0.6573 & 0.7824 & 200.4 ms \\
\quad no color seed & 0.4609 & 0.6052 & 186.2 ms \\
\quad no spatial prior & 0.6160 & 0.7476 & 200.6 ms \\
\quad single box & 0.6568 & 0.7819 & 192.9 ms \\
\bottomrule
\end{tabular}
\end{table}

The relative ranking on ADE20K is consistent with VOC: GrabCut $>$ Robust Superpixel $>$ Center Color, and the color seed remains the largest contributor. The absolute IoU values are moderately higher on ADE20K, likely because the Hugging Face mirror uses lower-resolution images (256--500~px) that reduce boundary-length-to-area ratios. The single-box ablation differs from the multi-box default by only 0.0005~IoU, replicating the VOC finding that box consensus provides essentially no average benefit. This preliminary cross-dataset evidence supports the transferability of the component analysis, though sample sizes and image resolutions differ materially. Full ADE20K validation (2,000 images) is deferred to future work.

\section{Discussion}

Representative evaluation changes the original interpretation. The proposed CPU method is reliably better than its minimal baseline, but not GrabCut; its value lies in transparency and a speed--accuracy operating point. Box localization remains SAM's principal signal, consistent with prior prompting studies \cite{gaus} and with BREPS's systematic finding of high sensitivity to bounding-box perturbations \cite{breps}; our added evidence is that channel sensitivity persists after realized prompt quality is approximately matched. SAM score selection becomes statistically reliable at full scale, although the improvement is only 0.0193 IoU and requires extra decoder calls.

Stratified analysis reveals systematic weaknesses: bicycles and other thin-structured objects consistently fail (IoU $<$ 0.21) because superpixel segmentation cannot capture wire-like geometry, and small targets (lowest area quartile) suffer across all CPU methods. The colour-seed component is the most area-sensitive ($+0.25$ IoU Q4$-$Q1), confirming that colour prototypes degrade on small regions. These per-class and per-size patterns---not just aggregate scores---should guide users in selecting the appropriate method for their target-object characteristics.

Negative results are equally informative. Multi-box consensus adds complexity without aggregate benefit. The strongest CPU method is classical GrabCut. On the hardest 20 samples, all CPU methods fail, but Center Color (0.14) and GrabCut (0.30) are more robust than the superpixel-based method (0.07), suggesting that structural priors can become liabilities when the underlying image segmentation is unreliable. These findings prevent a polished engineering artifact from being presented as a stronger methodological contribution than the evidence supports.

\paragraph{Practical guidance.}
Table~\ref{tab:guidance} distills the experimental evidence into actionable recommendations. The choice of method should depend on the target-object profile and available resources, not on aggregate benchmark ranking alone.

\begin{table}[ht]
\centering
\caption{Method selection guide based on target characteristics and constraints.}
\label{tab:guidance}
\begin{tabular}{p{0.28\linewidth}p{0.62\linewidth}}
\toprule
Scenario & Recommended method (and why) \\
\midrule
GPU available, box prompt feasible & SAM box-only (0.848 IoU; box dominates; point adds $<$0.002) \\
GPU available, point-only only & SAM point-only for animals/birds (0.72); Center Color for rigid objects (0.67 vs.\ SAM's 0.52 on tvmonitor) \\
CPU only, accuracy priority & GrabCut point+box (0.686 IoU); 8/1449 failures on large multi-textured objects---fall back to Center Color on error \\
CPU only, speed priority & Center Color (13 ms, 0.540 IoU) or Adaptive Superpixel (199 ms, 0.612 IoU) \\
Thin/wire-like targets (bicycles, chairs) & SAM box-only is the only method exceeding 0.30 IoU; all CPU methods fail ($<$0.21) \\
Large rigid objects (tv, sofa, bus) & GrabCut or SAM box-only; even Center Color reaches 0.67--0.67 \\
Small targets ($<12$K px) & Adaptive Superpixel (+0.014 over fixed); avoid GrabCut colour-seed ablation (0.28 IoU) \\
\bottomrule
\end{tabular}
\end{table}

\section{Threats to Validity}

Ground-truth-derived prompts are more controlled than real user input. Numeric matching of hit rate and box IoU does not make the metrics psychologically equivalent, and neither distribution is fitted to human errors. The benchmark selects one largest component per image rather than every instance. Only VOC 2012 and SAM ViT-B are tested at scale; ADE20K trials are preliminary (20 samples) and image resolution is lower. Full cross-dataset confirmation, SAM 2, larger checkpoints, and learned interactive baselines such as RITM or SimpleClick would improve external validity. GrabCut has eight explicit failures on VOC. The tuning split is class-balanced while validation has the natural distribution. Although source hashes prevent silent post-confirmatory changes, the hypotheses were motivated by exploratory results; an independent second dataset at scale would offer stronger confirmation.

\section{Conclusion}

Under a frozen full-validation protocol with 1,449 VOC images, we find that (1)~Robust Superpixel improves over Center Color but trails GrabCut; (2)~an adaptive superpixel-resolution variant yields a small, costless gain concentrated on small targets ($+0.014$ IoU in the lowest area quartile); (3)~SAM's box channel dominates, with point-only underperforming CPU baselines on several rigid-object classes; (4)~matched-quality box noise causes substantially larger loss than point noise ($+0.107$ IoU delta); and (5)~SAM score selection offers a modest but Holm-significant improvement over a single noisy prompt. Per-class and per-area stratification reveals that bicycles ($<$0.21 IoU for all CPU methods) and small targets are systematic failure modes, while GrabCut's eight failures are counter-intuitively concentrated on large multi-textured objects. A practical method-selection guide is provided. The primary contribution is a reproducible, stratified account of prompt-channel sensitivity, component-level evidence strength, and practical method trade-offs rather than a state-of-the-art claim.

\begin{thebibliography}{12}
\bibitem{voc} M. Everingham et al. The PASCAL Visual Object Classes Challenge. \emph{IJCV}, 2010. \url{https://www.robots.ox.ac.uk/~vgg/projects/pascal/VOC/}
\bibitem{grabcut} C. Rother, V. Kolmogorov, and A. Blake. GrabCut: Interactive Foreground Extraction Using Iterated Graph Cuts. \emph{ACM TOG}, 2004.
\bibitem{slic} R. Achanta et al. SLIC Superpixels Compared to State-of-the-Art Superpixel Methods. \emph{TPAMI}, 2012.
\bibitem{dextr} K.-K. Maninis et al. Deep Extreme Cut: From Extreme Points to Object Segmentation. \emph{CVPR}, 2018.
\bibitem{ritm} K. Sofiiuk, I. Petrov, and A. Konushin. Reviving Iterative Training with Mask Guidance for Interactive Segmentation. \emph{ICIP}, 2022.
\bibitem{simpleclick} Q. Liu et al. SimpleClick: Interactive Image Segmentation with Simple Vision Transformers. \emph{ICCV}, 2023.
\bibitem{sam} A. Kirillov et al. Segment Anything. \emph{ICCV}, 2023.
\bibitem{gaus} Y. Gaus et al. Performance Evaluation of Segment Anything Model with Variational Prompting. \emph{CVPRW}, 2024.
\bibitem{ppsam} A. Rahman et al. PP-SAM: Perturbed Prompts for Robust Adaptation of Segment Anything Model. \emph{CVPRW}, 2024.
\bibitem{uncertainsam} T. Kaiser et al. UncertainSAM: Segment Anything with Uncertainty Estimation. \emph{ICML}, 2025.
\bibitem{benenson} R. Benenson, S. Popov, and V. Ferrari. Large-Scale Interactive Object Segmentation with Human Annotators. \emph{CVPR}, 2019.
\bibitem{attackdefense} Liu et al. Attack for Defense: Adversarial Agents for Point Prompt Optimization. \emph{CVPR}, 2026.
\bibitem{ade20k} B. Zhou et al. Scene Parsing through ADE20K Dataset. \emph{CVPR}, 2017. \url{https://groups.csail.mit.edu/vision/datasets/ADE20K/}
\bibitem{sam2} N. Ravi et al. SAM 2: Segment Anything in Images and Videos. \emph{arXiv:2408.00714}, 2024.
\bibitem{efficientsam} Y. Xiong et al. EfficientSAM: Leveraged Masked Image Pretraining for Efficient Segment Anything. \emph{CVPR}, 2024.
\bibitem{hqsam} L. Ke et al. Segment Anything in High Quality. \emph{NeurIPS}, 2023.
\bibitem{breps} A. Moskalenko et al. BREPS: Bounding-Box Robustness Evaluation of Promptable Segmentation. \emph{AAAI}, 2026.
\end{thebibliography}

\end{document}
